\chapter{TỔNG KẾT}

\section{Kết quả đạt được}

Đồ án đã hoàn thành mục tiêu khảo sát, phân tích và thiết kế kiến trúc
tổng thể cho một hệ thống quản lý chuỗi nhà trọ ở quy mô doanh nghiệp
vừa. Các kết quả cụ thể bao gồm:

\begin{itemize}
\item
  Xác định rõ 04 tác nhân và 39 ca sử dụng (đánh số liên tục UC01--UC39,
  bao gồm UC39 Chatbot AI mới bổ sung) được tổ chức thành 6 nhóm chức
  năng chính, đảm bảo bao phủ toàn bộ nghiệp vụ vận hành. Đã triển khai
  backend thực tế đạt 38/39 UC kết nối API thật xuống MongoDB; chỉ còn
  UC17 ký số ở mức giả lập có chủ đích (xác nhận đổi trạng thái, chưa
  nhúng chữ ký CA). UC26 thanh toán online đã tích hợp thật cổng VNPay
  Sandbox (URL ký HMAC-SHA512, callback vnpay-return + IPN vnpay-ipn);
  chuyển khoản VietQR và tiền mặt đều áp dụng cơ chế đối soát 2 bước
  (pending\_cash --- Quản lý xác nhận rồi mới ghi nhận paid). UC35 xuất
  báo cáo đã sinh PDF thật bằng pdfkit (font DejaVu Sans tiếng Việt) qua
  GET /api/reports/pdf cùng xuất CSV UTF-8; xuất Excel .xlsx bằng
  exceljs để dành cho v2. UC36 gửi thông báo đã hoạt động thật qua email
  Gmail SMTP và thông báo in-app; kênh Telegram để dành cho v2.
\item
  Mô hình hoá đầy đủ 16 sơ đồ UML quan trọng (1 sơ đồ tổng quát, 5 sơ đồ
  Use Case con, 5 sơ đồ hoạt động, 1 sơ đồ lớp, 4 sơ đồ tuần tự).
\item
  Xác định 14 lớp đối tượng cốt lõi với thuộc tính, phương thức và mối
  quan hệ; tổ chức theo 5 gói (package) tương ứng với các ngữ cảnh
  nghiệp vụ. Khi triển khai thực tế trên MongoDB NoSQL, mô hình đã được
  tối ưu thành 11 collection độc lập (bổ sung Expense cho chi phí vận
  hành) và 4 lớp embedded (TaiSan trong Room, ChiTietHoaDon trong
  Invoice, ThongTinKhachThue và Otp trong User). Sơ đồ PlantUML phiên
  bản NoSQL có sẵn tại docs/class\_diagram.puml.
\item
  Đề xuất kiến trúc phân vai trò rõ ràng (Chủ trọ kiêm Admin, Quản lý,
  Khách thuê, Khách vãng lai), phù hợp các quy mô từ 1 đến hàng trăm nhà
  trọ.
\item
  Đề xuất phương án tích hợp cổng thanh toán (VNPay/MoMo/QR) và ký số
  trực tuyến - vốn là các "pain-point" lớn của quy trình thủ công.
\item
  Bộ mã nguồn PlantUML của các sơ đồ được lưu riêng trong thư mục
  docs/diagrams, giúp dễ dàng tái sinh sơ đồ và chuyển giao cho các nhóm
  phát triển, bảo trì sau này.
\end{itemize}

\section{Ưu điểm và khuyết điểm của hệ thống}

\subsection{Ưu điểm của hệ thống}

\begin{itemize}
\item
  Quản trị tập trung đa cơ sở: chủ trọ có cái nhìn tổng thể về toàn
  chuỗi, tránh tình trạng phân mảnh dữ liệu.
\item
  Phân quyền chi tiết theo vai trò (RBAC) giúp chủ trọ uỷ quyền vận hành
  cho quản lý mà vẫn kiểm soát được rủi ro - chính chủ trọ kiêm vai trò
  quản trị viên hệ thống và kế toán nên không bị phụ thuộc vào bên thứ
  ba để cấu hình nền tảng.
\item
  Tự động hoá tối đa các tác vụ lặp lại: tính tiền, sinh hoá đơn, gửi
  nhắc nợ.
\item
  Tích hợp ký số và cổng thanh toán điện tử, giảm thời gian thao tác và
  sai sót tiền mặt.
\item
  Trải nghiệm khách thuê được nâng cao thông qua cổng thông tin riêng,
  cho phép tra cứu minh bạch hợp đồng và hoá đơn, thanh toán online
  nhanh chóng.
\item
  Khả năng mở rộng nhờ kiến trúc multi-tenant và REST API, sẵn sàng tích
  hợp IoT (công tơ điện thông minh), AI dự báo lấp đầy hoặc các kênh
  thanh toán mới.
\end{itemize}

\subsection{Khuyết điểm của hệ thống}

\begin{itemize}
\item
  Phụ thuộc vào hạ tầng mạng và năng lực thiết bị di động của khách thuê
  - các khách thuê lớn tuổi có thể gặp khó khăn khi sử dụng cổng thông
  tin online.
\item
  Chi phí ban đầu cho chữ ký số và cổng thanh toán có thể cao so với nhà
  trọ quy mô rất nhỏ (dưới 5 phòng).
\item
  Việc nhập chỉ số điện nước thủ công vẫn cần thực hiện trên ứng dụng
  quản lý nếu chưa tích hợp IoT, dù đã giảm thời gian so với sổ tay.
\item
  Cần đầu tư đào tạo người dùng cuối (chủ trọ, quản lý) để sử dụng hết
  tính năng - đặc biệt là quản lý đa cơ sở.
\item
  Cần đầu tư về bảo mật và pháp lý (Nghị định 13/2023, hoá đơn điện tử)
  để tuân thủ quy định khi quy mô lớn.
\end{itemize}

\section{Hướng phát triển hệ thống trong tương lai}

\begin{itemize}
\item
  Tích hợp IoT đo lường: lắp đặt công tơ điện/nước thông minh truyền dữ
  liệu trực tiếp về hệ thống, loại bỏ hoàn toàn khâu nhập chỉ số thủ
  công.
\item
  Tích hợp AI dự báo: dự báo lấp đầy, xác định "khách có khả năng rời
  đi" để chủ trọ chủ động giữ chân; phát hiện gian lận thanh toán.
\item
  Xây dựng ứng dụng di động cho từng vai trò: app Khách thuê và app Quản
  lý trên cả Android và iOS, hỗ trợ chốt số điện nước và xác nhận thu
  tiền ngay tại hiện trường.
\item
  Mở rộng cho các loại hình bất động sản khác: căn hộ dịch vụ, homestay
  ngắn ngày, ký túc xá sinh viên - bằng cách thêm module và biểu mẫu hợp
  đồng phù hợp.
\item
  Triển khai mô hình SaaS đa khách hàng (multi-tenant) thương mại cho
  nhiều chủ trọ thuê bao theo gói tháng/năm.
\item
  Hỗ trợ marketplace tìm phòng: kết nối khách vãng lai và chủ trọ, hỗ
  trợ đặt cọc, đánh giá, xếp hạng cộng đồng.
\item
  Tích hợp eKYC (định danh điện tử) khi tạo hồ sơ khách thuê để chống
  giả mạo CCCD và tăng tính pháp lý của hợp đồng.
\item
  UC34 - Nâng cao báo cáo chi phí vận hành: collection Expense và biểu
  đồ chi phí đã có ở v1; v2 bổ sung endpoint GET
  /api/reports/profit-loss tổng hợp lãi/lỗ theo từng cơ sở và theo kỳ.
\item
  UC35 (nâng cấp v2) - Hoàn thiện xuất Excel/PDF báo cáo: v1 đã xuất PDF
  thật bằng pdfkit ở backend (GET /api/reports/pdf, kèm PDF hóa đơn GET
  /api/invoices/:id/pdf và PDF hợp đồng GET /api/contracts/:id/pdf,
  nhúng font DejaVu Sans tiếng Việt) và xuất CSV UTF-8 mở bằng Excel; v2
  bổ sung thư viện exceljs sinh file .xlsx mẫu doanh nghiệp và
  header/footer + watermark nâng cao cho PDF.
\item
  UC26 (nâng cấp v2) - Đưa cổng thanh toán lên môi trường production: v1
  đã tích hợp đầy đủ VNPay Sandbox --- sinh URL thanh toán ký
  HMAC-SHA512, callback GET /api/payments/vnpay-return, IPN callback GET
  /api/payments/vnpay-ipn, trang đối soát biên lai PaymentReturnPage; v2
  chuyển sang merchant VNPay production, bổ sung ví MoMo, biên lai PDF
  tự động và đối soát giao dịch định kỳ.
\item
  UC17 (nâng cấp v2) - Ký số thật bằng nhà cung cấp CA: nâng cấp
  endpoint PATCH /api/contracts/:id/sign hiện tại để nhúng chữ ký vẽ tay
  (react-signature-canvas) hoặc chữ ký số CA (VNPT-CA, FPT-CA, EasyCA)
  vào file PDF hợp đồng; tự động sinh PDF từ template Word.
\item
  UC36 (nâng cấp v2) - Mở rộng kênh thông báo: bổ sung Telegram Bot API
  bên cạnh email Gmail SMTP hiện có; áp dụng node-cron 8h sáng hằng ngày
  để tự động phát hành nhắc nợ cho các hoá đơn quá hạn.
\item
  UC23 (nâng cấp v2) - Tính tiền theo bậc thang EVN: nâng cấp logic POST
  /api/invoices/generate để hỗ trợ đơn giá điện theo bậc thang EVN 6 mức
  (1806đ/kWh đến 3157đ/kWh) bằng cách bổ sung field
  Service.tiers{[}\{from, to, price\}{]}; tương tự cho nước theo bậc
  thang công ty cấp nước địa phương.
\item
  Bổ sung Audit Log + Phân quyền chi tiết (RBAC theo permission): tạo
  collection AuditLog ghi nhận mọi thao tác nhạy cảm (login, đổi
  password, sửa hoá đơn đã thanh toán, xoá phòng\ldots); chuyển mô hình
  phân quyền từ role-based đơn giản sang permission-based để mỗi quản lý
  có thể được uỷ quyền chính xác nhóm thao tác nhất định.
\item
  Hỗ trợ song ngữ Việt-Anh (i18n bằng react-i18next), bổ sung bản đồ
  Việt Nam (leaflet/mapbox) hiển thị vị trí các cơ sở + sức khoẻ
  (xanh/vàng/đỏ theo doanh thu \& công nợ), và viết lại bộ test bằng
  Jest + Supertest cho API + Cypress cho E2E để thay thế bộ test hiện
  đang dựa trên mock class.
\end{itemize}

